\section{大数定律}

\begin{definition}
	设$\{f_n\}$是概率空间$(X,\mathscr{F},P)$上的随机变量序列。记$\mathscr{A}_n=\sigma(\{f_k:k\geqslant n\})$，由\cref{prop:SigmaField}(5)可知：
	\begin{equation*}
		\mathscr{A}=\underset{n=1}{\overset{+\infty}{\cap}}\mathscr{A}_n
	\end{equation*}
	也是一个$\sigma$域，称$\mathscr{A}$为$\{f_n\}$的\textbf{尾$\sigma$域}，$\mathscr{A}$中的事件被称为\textbf{尾事件}，称关于$\mathscr{A}$可测的随机变量为\textbf{尾随机变量}。
\end{definition}
\begin{note}
	从尾事件和尾随机变量的定义可以看出：一个事件如果是否发生只取决于“从某一项以后一直怎么样”，而不取决于前面有限项具体取什么值，那它就是尾事件。同理，一个随机变量如果它的取值只由“无穷远处的尾部”决定，而前面有限项怎么改都不影响它，那它就是尾随机变量。
\end{note}
\begin{theorem}[Kolmogorov 0-1 Law]\label{theo:01Law}
	设$\{f_n\}$是概率空间$(X,\mathscr{F},P)$上相互独立的随机变量序列，则：
	\begin{enumerate}
		\item 关于$\{f_n\}$的任一尾事件$A$，$P(A)$只能是$0$或者$1$；
		\item 关于$\{f_n\}$的任一尾随机变量为常数a.s.于$(X,\mathscr{A},P)$。
	\end{enumerate}
\end{theorem}
\begin{proof}
	(1)对任意的$n\in\mathbb{N}^+$，令$\mathscr{B}_n=\sigma(\{f_1,f_2,\dots,f_n\})$，由\cref{prop:Independent}(6)可知$\mathscr{B}_n$和$\mathscr{A}_{n+1}$独立，而$\mathscr{A}\subseteq\mathscr{A}_{n+1}$，所以$\mathscr{B}_n$和$\mathscr{A}$独立，$\underset{n=1}{\overset{+\infty}{\cup}}\mathscr{B}_n$与$\mathscr{A}$独立。令$\mathscr{B}=\sigma\left(\underset{n=1}{\overset{+\infty}{\bigcup}}\mathscr{B}_n\right)$，由$\mathscr{B}_n$的递增性和\cref{prop:SigmaField}(2)可知$\underset{n=1}{\overset{+\infty}{\cup}}\mathscr{B}_n$是$\pi$系，根据\cref{prop:Independent}(2)可得$\mathscr{B}$和$\mathscr{A}$独立。因为$f_n$关于$\mathscr{B}_n$可测，所以$\sigma(\{f_n\})\subseteq\mathscr{B}$；因为$\mathscr{B}_n\subseteq\sigma(\{f_n\})$，所以$\mathscr{B}\subseteq\sigma(\{f_n\})$，即$\mathscr{B}=\sigma(\{f_n\})$。任取$A\in\mathscr{A}$，则$A\in\mathscr{A}_1=\sigma(\{f_n\})$，即$A\in\mathscr{B}$。将$A$表示为$A\cap A$，由$\mathscr{A}$和$\mathscr{B}$的独立性可得：
	\begin{equation*}
		P(A)=P(A\cap A)=[P(A)]^2
	\end{equation*}
	所以$P(A)$为$0$或$1$。\par
	(2)任取关于$\{f_n\}$的尾随机变量$f$，由\cref{prop:MeasurableFunction}(1.b)可知对任意的$a\in\mathbb{R}^{}$，$\{f\leqslant a\}$是尾事件，根据(1)可知$P(f\leqslant a)$为$0$或$1$。取：
	\begin{equation*}
		a_0=\inf\{a\in\mathbb{R}^{}:P(f\leqslant a)=1\}
	\end{equation*}
	由分布函数的右连续性可得$P(f\leqslant a_0)=1$。根据\cref{prop:Measure}(3)（下连续性）和\cref{prop:RSeq}(5)可得：
	\begin{equation*}
		P(f<a_0)=\lim_{n\to+\infty}P\left(f\leqslant a_0-\frac{1}{n}\right)=0
	\end{equation*}
	由\cref{prop:Measure}(2)可得$P(f=a_0)=1$。
\end{proof}
\begin{note}
	Kolmogorov 0-1 Law说明了如果$\{f_n\}$是独立的随机变量序列，则：
	\begin{enumerate}
		\item $\sum\limits_{n=1}^{+\infty}f_n$要么收敛a.s.于$(X,\mathscr{A},P)$要么发散a.s.于$(X,\mathscr{A},P)$；
		\item 若$\lim\limits_{n\to+\infty}\dfrac{1}{n}\sum\limits_{i=1}^{n}f_i$存在，则只能是一个常数a.s.于$(X,\mathscr{A},P)$。
	\end{enumerate}
	由函数项级数收敛的Cauchy式定义、\cref{prop:SigmaField}(2)和\cref{prop:MeasurableFunction}(5.b)(3)可以得到$\sum\limits_{n=1}^{+\infty}f_n$收敛是尾事件，根据\cref{prop:MeasurableFunction}(5.a)(6)和\cref{prop:RSeq}(8.b)可知如果$\lim\limits_{n\to+\infty}\dfrac{1}{n}\sum\limits_{i=1}^{n}f_i$存在，则$\lim\limits_{n\to+\infty}\dfrac{1}{n}\sum\limits_{i=1}^{n}f_i$是尾随机变量。
\end{note}
\begin{lemma}[Borel-Cantelli Lemma]\label{lem:BorelCantelli}
	设$(X,\mathscr{F},P)$是概率空间，$\{A_n\}\subseteq\mathscr{F}$。
	\begin{enumerate}
		\item 若$\sum\limits_{n=1}^{+\infty}P(A_n)<+\infty$，则：
		\begin{equation*}
			P\left(\varlimsup_{n\to+\infty}A_n\right)=0
		\end{equation*}
		\item 若$\{A_n\}$相互独立且$\sum\limits_{n=1}^{+\infty}P(A_n)=+\infty$，则：
		\begin{equation*}
			P\left(\varlimsup_{n\to+\infty}A_n\right)=1
		\end{equation*}
	\end{enumerate}
\end{lemma}
\begin{proof}
	(1)由\cref{prop:SetLimit}(3)、\cref{prop:Measure}(3)（上连续性，次可列可加性）和\cref{prop:RSeq}(6)可得：
	\begin{equation*}
		P\left(\varlimsup_{n\to+\infty}A_n\right)=P\left(\underset{n=1}{\overset{+\infty}{\cap}}\underset{k=n}{\overset{+\infty}{\cup}}A_k\right)=\lim_{n\to+\infty}P\left(\underset{k=n}{\overset{+\infty}{\cup}}A_k\right)\leqslant\lim_{n\to+\infty}\sum_{k=n}^{+\infty}P(A_k)=0
	\end{equation*}\par
	(2)由\cref{prop:SetOperation}(7)、\cref{prop:SetLimit}(3)、\cref{prop:Measure}(3)（下连续性，上连续性）和\cref{prop:RSeq}(6)可得：
	\begin{align*}
		&P\left[\left(\varlimsup_{n\to+\infty}A_n\right)^c\right]=P\left(\underset{n=1}{\overset{+\infty}{\cup}}\underset{k=n}{\overset{+\infty}{\cap}}A_k^c\right)=\lim_{n\to+\infty}P\left(\underset{k=n}{\overset{+\infty}{\cap}}A_k^c\right)=\lim_{n\to+\infty}\lim_{m\to+\infty}P\left(\underset{k=n}{\overset{m}{\cap}}A_k^c\right) \\
		=&\lim_{n\to+\infty}\lim_{m\to+\infty}\prod_{k=n}^{m}[1-P(A_k)]=\lim_{n\to+\infty}\lim_{m\to+\infty}\exp\left\{\sum_{k=n}^{m}\ln[1-P(A_k)]\right\} \\
		\leqslant&\lim_{n\to+\infty}\lim_{m\to+\infty}\exp\left\{-\sum_{k=n}^{m}P(A_k)\right\}=\lim_{n\to+\infty}\exp\left\{-\sum_{k=n}^{+\infty}P(A_k)\right\}=0
	\end{align*}
	由\cref{prop:Measure}(2)即可得出结论。
\end{proof}

\subsection{弱大数定律}
\begin{definition}
	设$\{X_n\}$是一个随机变量序列，若
	\begin{equation*}
		\frac{1}{n}\sum_{i=1}^{n}X_i\overset{P}{\longrightarrow}\frac{1}{n}\sum_{i=1}^{n}\operatorname{E}(X_i)
	\end{equation*}
	则称$\{X_n\}$服从\gls{WeakLawOfLargeNumbers}。
\end{definition}
\begin{theorem}[Markov's Weak Law of Large Numbers]
	设$\{X_n\}$是一个随机变量序列，若：
	\begin{equation*}
		\lim_{n\to+\infty}\left[\frac{1}{n^2}\operatorname{Var}\left(\sum_{i=1}^{n}X_i\right)\right]=0
	\end{equation*}
	则$\{X_n\}$服从弱大数定律。
\end{theorem}
\begin{proof}
	由\cref{ineq:Chebyshev}可知对任意的$\varepsilon>0$有：
	\begin{equation*}
		P\left(\left|\frac{1}{n}\sum_{i=1}^{n}X_i-\frac{1}{n}\sum_{i=1}^{n}\operatorname{E}(X_i)\right|\geqslant\varepsilon\right)\leqslant\frac{\operatorname{Var}\left(\dfrac{1}{n}\sum\limits_{i=1}^{n}X_i\right)}{\varepsilon^2}=\frac{1}{n^2\varepsilon^2}\operatorname{Var}\left(\sum_{i=1}^{n}X_i\right)\to0\qedhere
	\end{equation*}
\end{proof}
\begin{corollary}
	由Markov's Weak Law of Large Numbers可以推出：
	\begin{enumerate}
		\item (Chebyshev's Weak Law of Large Numbers)设$\{X_n\}$为一列互不相关的随机变量，若对于任意的$n\in\mathbb{N}^+$，有$\operatorname{Var}(X_i)\leqslant c<+\infty$，则$\{X_n\}$服从弱大数定律。
		\item (Bernoulli's Weak Law of Large Numbers)设$\{X_n\}$为一列独立同两点分布$b(1,p)$的随机变量，则$\{X_n\}$服从弱大数定律。
	\end{enumerate}
\end{corollary}
\begin{proof}
	(1)由\cref{prop:MeasurableIntegral}(6)、\cref{prop:Variance}(3)和\cref{prop:CovMat}(7)可得：
	\begin{equation*}
		\operatorname{Var}\left(\frac{1}{n}\sum_{i=1}^{n}X_i\right)=\frac{1}{n^2}\sum_{i=1}^{n}\operatorname{Var}(X_i)\leqslant\frac{1}{n^2}nc=\frac{c}{n}\to0
	\end{equation*}\par
	(2)由两点分布的性质或者(1)直接可得。
\end{proof}
%\begin{theorem}[khinchin's Weak Law of Large Numbers]\label{theo:WeakLawOfLargeNumbers}
%%	设$\{Xm_n\}$是一列独立同分布的随机变量，若$X_n$的数学期望有限，则$\{X_n\}$服从大数定律。
%\end{theorem}
\begin{proof}
	\info{需要写完特征函数}
\end{proof}

\begin{theorem}[Strong Law of Large Numbers]
	\label{theo:StrongLawOfLargeNumbers}
	设$\{X_n\}$为一列独立同分布的随机变量，若$\operatorname{E}(X_1)$有意义，则有:
	\begin{equation*}
		P\left(\left\{\lim_{n\to+\infty}\bar{X}_n\ne\operatorname{E}(X_1)\right\}\right)=0
	\end{equation*}
	即$\bar{X}_n\overset{\text{a.e.}}{\longrightarrow}\operatorname{E}(X_1)$。
\end{theorem}
